\documentclass[10pt,a4paper]{amsart}
\usepackage{amsmath,amssymb,amsthm,graphicx}
\usepackage{verbatim}
\newenvironment{reference}{\comment}{\endcomment}
%\newenvironment{reference}{}{}
\newenvironment{slogan}{\comment}{\endcomment}
\newenvironment{history}{\comment}{\endcomment}
% For commutative diagrams we use Xy-pic
\usepackage[all]{xy}

% We use 2cell for 2-commutative diagrams.
\xyoption{2cell}
% Theorem environments.
%
\theoremstyle{plain}
\newtheorem{theorem}[subsection]{Theorem}
\newtheorem{proposition}[subsection]{Proposition}
\newtheorem{lemma}[subsection]{Lemma}

\theoremstyle{definition}
\newtheorem{definition}[subsection]{Definition}
\newtheorem{example}[subsection]{Example}
\newtheorem{exercise}[subsection]{Exercise}
\newtheorem{situation}[subsection]{Situation}

\theoremstyle{remark}
\newtheorem{remark}[subsection]{Remark}
\newtheorem{remarks}[subsection]{Remarks}

\numberwithin{equation}{subsection}

% Macros
%
\def\lim{\mathop{\mathrm{lim}}\nolimits}
\def\colim{\mathop{\mathrm{colim}}\nolimits}
\def\Spec{\mathop{\mathrm{Spec}}}
\def\Hom{\mathop{\mathrm{Hom}}\nolimits}
\def\Ext{\mathop{\mathrm{Ext}}\nolimits}
\def\SheafHom{\mathop{\mathcal{H}\!\mathit{om}}\nolimits}
\def\SheafExt{\mathop{\mathcal{E}\!\mathit{xt}}\nolimits}
\def\Sch{\mathit{Sch}}
\def\Mor{\mathop{\mathrm{Mor}}\nolimits}
\def\Ob{\mathop{\mathrm{Ob}}\nolimits}
\def\Sh{\mathop{\mathit{Sh}}\nolimits}
\def\NL{\mathop{N\!L}\nolimits}
\def\CH{\mathop{\mathrm{CH}}\nolimits}
\def\proetale{{pro\text{-}\acute{e}tale}}
\def\etale{{\acute{e}tale}}
\def\QCoh{\mathit{QCoh}}
\def\Ker{\mathop{\mathrm{Ker}}}
\def\Im{\mathop{\mathrm{Im}}}
\def\Coker{\mathop{\mathrm{Coker}}}
\def\Coim{\mathop{\mathrm{Coim}}}

% Boxtimes
%
\DeclareMathSymbol{\boxtimes}{\mathbin}{AMSa}{"02}

%
% Macros for moduli stacks/spaces
%
\def\QCohstack{\mathcal{QC}\!\mathit{oh}}
\def\Cohstack{\mathcal{C}\!\mathit{oh}}
\def\Spacesstack{\mathcal{S}\!\mathit{paces}}
\def\Quotfunctor{\mathrm{Quot}}
\def\Hilbfunctor{\mathrm{Hilb}}
\def\Curvesstack{\mathcal{C}\!\mathit{urves}}
\def\Polarizedstack{\mathcal{P}\!\mathit{olarized}}
\def\Complexesstack{\mathcal{C}\!\mathit{omplexes}}
% \Pic is the operator that assigns to X its picard group, usage \Pic(X)
% \Picardstack_{X/B} denotes the Picard stack of X over B
% \Picardfunctor_{X/B} denotes the Picard functor of X over B
\def\Pic{\mathop{\mathrm{Pic}}\nolimits}
\def\Picardstack{\mathcal{P}\!\mathit{ic}}
\def\Picardfunctor{\mathrm{Pic}}
\def\Deformationcategory{\mathcal{D}\!\mathit{ef}}
\usepackage{cleveref}
\crefname{originallemma}{original lemma}{original lemmas}
\Crefname{originallemma}{Original lemma}{Original lemmas}
\makeatletter
\newlabel{lemma-cd-local}{{lemma-cd-local}{source}{}{sourcepointer-1}{}}
\newlabel{lemma-cd-local@cref}{{[originallemma][0][]lemma-cd-local}{[1][0][]source}{}{}{}}
\makeatother
\crefname{originallemma}{original lemma}{original lemmas}
\Crefname{originallemma}{Original lemma}{Original lemmas}
\makeatletter
\newlabel{lemma-cd}{{lemma-cd}{source}{}{sourcepointer-2}{}}
\newlabel{lemma-cd@cref}{{[originallemma][0][]lemma-cd}{[1][0][]source}{}{}{}}
\makeatother
\crefname{originalreference}{original reference}{original references}
\Crefname{originalreference}{Original reference}{Original references}
\makeatletter
\newlabel{algebra-lemma-spec-spectral}{{algebra-lemma-spec-spectral}{source}{}{sourcepointer-3}{}}
\newlabel{algebra-lemma-spec-spectral@cref}{{[originalreference][0][]algebra-lemma-spec-spectral}{[1][0][]source}{}{}{}}
\makeatother
\crefname{originalreference}{original reference}{original references}
\Crefname{originalreference}{Original reference}{Original references}
\makeatletter
\newlabel{topology-lemma-spectral-sub}{{topology-lemma-spectral-sub}{source}{}{sourcepointer-4}{}}
\newlabel{topology-lemma-spectral-sub@cref}{{[originalreference][0][]topology-lemma-spectral-sub}{[1][0][]source}{}{}{}}
\makeatother
\crefname{originalreference}{original reference}{original references}
\Crefname{originalreference}{Original reference}{Original references}
\makeatletter
\newlabel{cohomology-proposition-cohomological-dimension-spectral}{{cohomology-proposition-cohomological-dimension-spectral}{source}{}{sourcepointer-5}{}}
\newlabel{cohomology-proposition-cohomological-dimension-spectral@cref}{{[originalreference][0][]cohomology-proposition-cohomological-dimension-spectral}{[1][0][]source}{}{}{}}
\makeatother
\crefname{originalreference}{original reference}{original references}
\Crefname{originalreference}{Original reference}{Original references}
\makeatletter
\newlabel{cohomology-proposition-vanishing-Noetherian}{{cohomology-proposition-vanishing-Noetherian}{source}{}{sourcepointer-6}{}}
\newlabel{cohomology-proposition-vanishing-Noetherian@cref}{{[originalreference][0][]cohomology-proposition-vanishing-Noetherian}{[1][0][]source}{}{}{}}
\makeatother
\crefname{originalreference}{original reference}{original references}
\Crefname{originalreference}{Original reference}{Original references}
\makeatletter
\newlabel{algebra-lemma-support-closed}{{algebra-lemma-support-closed}{source}{}{sourcepointer-7}{}}
\newlabel{algebra-lemma-support-closed@cref}{{[originalreference][0][]algebra-lemma-support-closed}{[1][0][]source}{}{}{}}
\makeatother
\crefname{originalreference}{original reference}{original references}
\Crefname{originalreference}{Original reference}{Original references}
\makeatletter
\newlabel{dualizing-lemma-local-cohomology-and-restriction}{{dualizing-lemma-local-cohomology-and-restriction}{source}{}{sourcepointer-8}{}}
\newlabel{dualizing-lemma-local-cohomology-and-restriction@cref}{{[originalreference][0][]dualizing-lemma-local-cohomology-and-restriction}{[1][0][]source}{}{}{}}
\makeatother
\title{Local cohomology of a finite module vanishes above the dimension of its support, and cohomological dimension is bounded by the dimension of the ring}
\author{The Stacks Project Authors}
\date{stacks-project:local-cohomology.tex:local-snapshot-20261012; GNU FDL 1.2 or later}
\begin{document}
\maketitle
\noindent\textit{Editorial note.} All delivered mathematical text is the author's own, copied from the pinned Stacks Project chapter source sources/stacks-local-cohomology/pinned-original-source.tex; no statement or proof step was rewritten, reordered or removed.\par\medskip
\noindent\textit{Editorial note.} Preamble provenance: the Stacks Project keeps one single shared preamble for every chapter; this card takes it directly from the staged sources/stacks-local-cohomology/preamble.tex as three verbatim spans (lines 8-12, 14-18 and 154-217), so no derived file was written and nothing was added or redefined. The wrapper class and the base packages are supplied by the assembler.\par\medskip
\noindent\textit{Editorial note.} Bibliography provenance: neither the counted statement nor its proof contains a citation (this lane rescans the whole span before writing the card), so the card carries no bibliography and no entry was reconstructed.\par\medskip
\noindent\textit{Editorial note.} Selection provenance: this card was opened under the file-based control contract (CONTROL.md), with the mechanical pre-card gate run on its statement and proof interval before the card was written.\par\medskip
\noindent\textit{Allowed standard prerequisites (editorial).}\begin{itemize}
\item The reduction of the cohomological dimension to the local case, and the characterisation of cohomological dimension by the vanishing of the cohomology of quasi-coherent sheaves on the complement. Both are invoked inside the counted proof (source lines 769 and 781) as pointers to items of this chapter. Reduce the second assertion to a local ring and translate the required vanishing into a statement about the complement.
\item The description of the spectrum of a quotient as the support of the corresponding module, and the identification of the local cohomology of a module over the quotient with the local cohomology of the ideal extended to it. Both are invoked inside the counted proof (source lines 789-794) as cross-chapter pointers. Reduce the first assertion to the second applied to the quotient by the annihilator.
\end{itemize}\medskip
\ifdefined\hypertarget\hypertarget{sourcepointer-1}{}\fi
\noindent\textit{Source reference.} Original lemma lemma-cd-local: sources/stacks-local-cohomology/pinned-original-source.tex, source lines 769--769. Navigation pointer only: the counted text refers to the Stacks item lemma-cd-local, which belongs to another chapter of the project and is not part of this card. The pointer carries the target's own label as its number because this chapter's staged snapshot does not contain the other chapter's tag table; its content is not reproduced here.\par\medskip
\ifdefined\hypertarget\hypertarget{sourcepointer-2}{}\fi
\noindent\textit{Source reference.} Original lemma lemma-cd: sources/stacks-local-cohomology/pinned-original-source.tex, source lines 679--679. Navigation pointer only: the counted text refers to the Stacks item lemma-cd, which belongs to another chapter of the project and is not part of this card. The pointer carries the target's own label as its number because this chapter's staged snapshot does not contain the other chapter's tag table; its content is not reproduced here.\par\medskip
\ifdefined\hypertarget\hypertarget{sourcepointer-3}{}\fi
\noindent\textit{Source reference.} Original reference algebra-lemma-spec-spectral: sources/stacks-local-cohomology/pinned-original-source.tex, source lines 776--776. Navigation pointer only: the counted text refers to the Stacks item algebra-lemma-spec-spectral, which belongs to another chapter of the project and is not part of this card. The pointer carries the target's own label as its number because this chapter's staged snapshot does not contain the other chapter's tag table; its content is not reproduced here.\par\medskip
\ifdefined\hypertarget\hypertarget{sourcepointer-4}{}\fi
\noindent\textit{Source reference.} Original reference topology-lemma-spectral-sub: sources/stacks-local-cohomology/pinned-original-source.tex, source lines 777--777. Navigation pointer only: the counted text refers to the Stacks item topology-lemma-spectral-sub, which belongs to another chapter of the project and is not part of this card. The pointer carries the target's own label as its number because this chapter's staged snapshot does not contain the other chapter's tag table; its content is not reproduced here.\par\medskip
\ifdefined\hypertarget\hypertarget{sourcepointer-5}{}\fi
\noindent\textit{Source reference.} Original reference cohomology-proposition-cohomological-dimension-spectral: sources/stacks-local-cohomology/pinned-original-source.tex, source lines 779--779. Navigation pointer only: the counted text refers to the Stacks item cohomology-proposition-cohomological-dimension-spectral, which belongs to another chapter of the project and is not part of this card. The pointer carries the target's own label as its number because this chapter's staged snapshot does not contain the other chapter's tag table; its content is not reproduced here.\par\medskip
\ifdefined\hypertarget\hypertarget{sourcepointer-6}{}\fi
\noindent\textit{Source reference.} Original reference cohomology-proposition-vanishing-Noetherian: sources/stacks-local-cohomology/pinned-original-source.tex, source lines 784--784. Navigation pointer only: the counted text refers to the Stacks item cohomology-proposition-vanishing-Noetherian, which belongs to another chapter of the project and is not part of this card. The pointer carries the target's own label as its number because this chapter's staged snapshot does not contain the other chapter's tag table; its content is not reproduced here.\par\medskip
\ifdefined\hypertarget\hypertarget{sourcepointer-7}{}\fi
\noindent\textit{Source reference.} Original reference algebra-lemma-support-closed: sources/stacks-local-cohomology/pinned-original-source.tex, source lines 790--790. Navigation pointer only: the counted text refers to the Stacks item algebra-lemma-support-closed, which belongs to another chapter of the project and is not part of this card. The pointer carries the target's own label as its number because this chapter's staged snapshot does not contain the other chapter's tag table; its content is not reproduced here.\par\medskip
\ifdefined\hypertarget\hypertarget{sourcepointer-8}{}\fi
\noindent\textit{Source reference.} Original reference dualizing-lemma-local-cohomology-and-restriction: sources/stacks-local-cohomology/pinned-original-source.tex, source lines 794--794. Navigation pointer only: the counted text refers to the Stacks item dualizing-lemma-local-cohomology-and-restriction, which belongs to another chapter of the project and is not part of this card. The pointer carries the target's own label as its number because this chapter's staged snapshot does not contain the other chapter's tag table; its content is not reproduced here.\par\medskip
\begin{lemma}
\label{lemma-cd-dimension}
Let $I \subset A$ be a finitely generated ideal of a ring $A$.
If $M$ is a finite $A$-module, then
$H^i_{V(I)}(M) = 0$ for $i > \dim(\text{Supp}(M))$.
In particular, we have $\text{cd}(A, I) \leq \dim(A)$.
\end{lemma}

\begin{proof}
We first prove the second statement.
Recall that $\dim(A)$ denotes the Krull dimension. By
Lemma \ref{lemma-cd-local} we may assume $A$ is local.
If $V(I) = \emptyset$, then the result is true.
If $V(I) \not = \emptyset$, then
$\dim(\Spec(A) \setminus V(I)) < \dim(A)$ because
the closed point is missing. Observe that
$U = \Spec(A) \setminus V(I)$ is a quasi-compact
open of the spectral space $\Spec(A)$, hence a spectral space itself.
See Algebra, Lemma \ref{algebra-lemma-spec-spectral} and
Topology, Lemma \ref{topology-lemma-spectral-sub}.
Thus Cohomology, Proposition
\ref{cohomology-proposition-cohomological-dimension-spectral}
implies $H^i(U, \mathcal{F}) = 0$ for $i \geq \dim(A)$
which implies what we want by Lemma \ref{lemma-cd}.
In the Noetherian case the reader may use
Grothendieck's Cohomology, Proposition
\ref{cohomology-proposition-vanishing-Noetherian}.

\medskip\noindent
We will deduce the first statement from the second.
Let $\mathfrak a$ be the annihilator of the finite $A$-module $M$.
Set $B = A/\mathfrak a$. Recall that $\Spec(B) = \text{Supp}(M)$, see
Algebra, Lemma \ref{algebra-lemma-support-closed}.
Set $J = IB$. Then $M$ is a $B$-module
and $H^i_{V(I)}(M) = H^i_{V(J)}(M)$, see
Dualizing Complexes, Lemma
\ref{dualizing-lemma-local-cohomology-and-restriction}.
Since $\text{cd}(B, J) \leq \dim(B) = \dim(\text{Supp}(M))$
by the first part we conclude.
\end{proof}

\end{document}
